\section{课程定位与 NLP 历史：Transformer 不是凭空出现的}

这堂 V4 开场课承担两个任务。第一，它给后续嘉宾讲座建立共同词汇：attention、large language model、alignment、reasoning 与 agent。第二，它展示一条尚未收束的研究路线：模型能力快速扩张，但数据效率、持续学习、记忆、可靠性和安全边界仍然缺口明显。阅读时应把 2024 年课堂中的产品例子与更稳定的机制分开，前者是时间点证据，后者才是可迁移的知识。

\lecturefigure{slide-001.jpg}{CS25 Transformers United V4 开场页}{官方 CS25 V4 slide deck，第 1 页}

\begin{knowledgebox}{读图：Overview 不是“快速过一遍名词”}
标题页列出四位课程组织者，后续内容也确实由不同讲者分段完成。它覆盖历史、机制、应用、弱点、推理与 agent，因此更像一张课程地图。讲义会补足公式和系统边界，但不会把课堂未出现的 dashboard、incident runbook 或合规流程塞回去。
\end{knowledgebox}

\teachervoice{开场说明 CS25 的目标不是重复教材，而是让正在做研究的人直接解释新结果，使学生能够把想法带回自己的研究或形成合作。这个定位决定了课堂会并列呈现共识、开放问题和讲者个人判断。}

\subsection{课程希望学生带走什么}

技术概览容易退化为模型名称列表。课程给出的学习目标更具体：理解 Transformer 怎样工作，理解它为何能跨越 NLP，识别前沿研究方向，并持续追问剩余的弱点。这个框架要求我们既解释成功条件，也解释失败边界。读下面的目标页时，先看四个动词分别指向 mechanism、application、research 与 limitation，而不是把它当作课程宣传页。

\lecturefigure{slide-011.jpg}{课程的四类学习目标}{官方 CS25 V4 slide deck，第 11 页}

\begin{importantbox}{本讲的四条主线}
\begin{enumerate}
  \item \textbf{Mechanism}：attention、self-attention、multi-head 与 cross-attention 怎样传递信息；
  \item \textbf{Scale and training}：参数、数据、alignment 与 preference optimization 怎样塑造能力；
  \item \textbf{Adaptation}：记忆、检索、持续学习、model editing 与 reasoning scaffold 怎样补模型缺口；
  \item \textbf{System}：agent、tool、memory、multi-agent communication 与 safety 怎样把模型变成可执行系统。
\end{enumerate}
\end{importantbox}

\subsection{从 rule system 到 attention timeline}

课程先给出时间轴，不是为了争论某项技术的“第一篇论文”，而是为了说明每一代方法在修补什么瓶颈。Rule-based NLP 把语言写成规则；word embedding 把离散词映成连续向量；RNN/LSTM 处理序列；attention 让模型对不同位置分配不同权重；Transformer 再把 attention 变成主体计算结构。

\lecturefigure{slide-013.jpg}{从规则系统到 Transformer 与多模态模型的时间轴}{官方 CS25 V4 slide deck，第 13 页}

\begin{knowledgebox}{术语消化：时间轴上的五个台阶}
\begin{tabularx}{\textwidth}{L{0.19\textwidth} X}
\toprule
阶段 & 它主要解决的问题 \\
\midrule
Rule-based system & 用人工规则编码句法与语义，透明但覆盖脆弱，迁移成本高。 \\
Word embedding & 用连续向量表达相似性，使模型能共享词之间的统计结构。 \\
RNN/LSTM & 用 recurrent state 处理可变长序列，但训练串行，长距离信息容易衰减。 \\
Attention & 对输入位置做内容相关加权，不再把全部历史压进单个状态。 \\
Transformer & 以 attention 和前馈网络为主干，扩大并行训练和规模化的空间。 \\
\bottomrule
\end{tabularx}
\end{knowledgebox}

\lecturefigure{slide-014.jpg}{早期 NLP 的歧义、规则维护与泛化难题}{官方 CS25 V4 slide deck，第 14 页}

\begin{warningbox}{历史页不支持“旧方法完全无用”}
规则、词典、检索、有限状态机和传统特征仍然适合高约束任务。Transformer 的优势在于端到端表示学习和规模化，不是让所有结构化方法失效。工程系统通常会把 learned model 与规则、数据库、验证器一起使用。
\end{warningbox}

\lecturefigure{slide-015.jpg}{NLP 方法从手写规则走向学习式表示}{官方 CS25 V4 slide deck，第 15 页}

\lecturefigure{slide-016.jpg}{ELIZA 展示 rule-based 对话系统的能力与边界}{官方 CS25 V4 slide deck，第 16 页}

\begin{knowledgebox}{读图：ELIZA 为什么能“像在对话”}
ELIZA 使用模式匹配与模板重写，把用户句子中的片段填入预设回答。它说明表面流畅不等于内部理解：只要交互形式足够合适，有限规则也能制造连贯感。今天评估 LLM 时，同样不能把语言自然度直接等同于事实性、推理正确性或行动可靠性。
\end{knowledgebox}

\lecturefigure{slide-017.jpg}{语言学规则、semantic parsing 与早期符号表示}{官方 CS25 V4 slide deck，第 17 页}

\begin{importantbox}{结构知识仍然有价值}
Semantic parsing 把自然语言映射成逻辑形式或可执行表示。现代 LLM 可以减少手写规则，却没有消除结构需求：tool calling 的 schema、SQL、程序、证明步骤和 agent protocol 都是在重新引入可验证结构。
\end{importantbox}

\lecturefigure{slide-018.jpg}{Word embedding 将离散词映射到连续向量空间}{官方 CS25 V4 slide deck，第 18 页}

\begin{importantbox}{Embedding 的最小定义}
设词表项 $w$ 的向量为 $e_w\in\mathbb{R}^{d}$。Embedding（嵌入）把离散符号映射到 $d$ 维连续空间，使相似词在某种训练目标下获得接近的表示。常见相似度是
\[
\operatorname{cos}(e_i,e_j)=\frac{e_i^\top e_j}{\lVert e_i\rVert_2\lVert e_j\rVert_2}.
\]
其中 $e_i,e_j$ 是两个 token 或对象的向量，分子是内积，分母消除向量长度影响。接近只表示表示空间中的相关性，不自动等于事实关系或因果关系。
\end{importantbox}

\subsection{本章小结}

Transformer 的历史位置可以压缩成一句话：它把“从序列中选择相关信息”提升为可并行、可堆叠、可扩展的主要计算模式。后续能力并非来自单一公式，而是 architecture、data、optimization、hardware 与 interface 共同演化的结果。

